% Build twice with pdflatex for the table of contents and cross-references.
% No external bibliography processor or downloaded figure is required.
\documentclass[11pt,a4paper]{article}
\usepackage[T1]{fontenc}
\usepackage[utf8]{inputenc}
\usepackage[margin=24mm]{geometry}
\usepackage{amsmath,amssymb,bm}
\usepackage{booktabs,longtable,array}
\usepackage{enumitem}
\usepackage[colorlinks=true,linkcolor=blue,citecolor=blue,urlcolor=blue]{hyperref}
\usepackage{listings}
\lstset{basicstyle=\ttfamily\small,breaklines=true,columns=fullflexible,frame=single}
\setlength{\emergencystretch}{2em}
\newcommand{\dd}{\,\mathrm{d}}
\newcommand{\tr}{^{\mathsf T}}
\newcommand{\R}{\mathbb R}
\newcommand{\Id}{\bm I}
\newcommand{\vn}{\bm n}
\newcommand{\vN}{\bm N_0}
\newcommand{\vt}{\bm t}
\newcommand{\vq}{\bm q}
\newcommand{\vs}{\bm s}
\newcommand{\vC}{\bm C}
\newcommand{\Proj}{\bm P}
\newcommand{\Nm}{\bm N}
\newcommand{\Tm}{\bm T}
\newcommand{\code}[1]{\texttt{\detokenize{#1}}}
\DeclareMathOperator{\diag}{diag}
\title{Proposed MORIS Implementation of Mlika Frictional Contact\\
\large Thesis Eq.~(4.25), reusing the corrected Seitz assembly}
\author{Implementation design for review}
\date{Revised November 25, 2026}
\begin{document}
\maketitle

\begin{abstract}
This document proposes a separate Mlika frictional-contact IWG implementing
Eq.~(4.25) of Mlika's thesis, equivalently Eq.~(15) of
Mlika--Renard--Chouly (2017). The implementation reuses the corrected Seitz
IWG's assembly organization: common ray mapping, current and previous field
interpolators, a quadrature-point Coulomb projection, separate leader/follower
residuals, current Jacobian blocks, and a previous-slab Jacobian for adjoints.
Both laws use current constitutive traction; previous configurations enter only
the slip kinematics. Mlika uses velocity, while Seitz uses an increment. A small
Mlika-specific adapter is proposed instead of duplicating the Seitz assembly.
No independent internal friction history or additional time-scale property is
introduced. This is a revised design proposal; Mlika implementation remains on
hold pending review. The dormant friction prototype and traction-history map
have since been removed from the existing Mlika class. That class retains only
the verified frictionless energy formulation, Eqs.~(4.28)--(4.30); frictional
contact uses the corrected Seitz implementation. The adapter described below
remains an unimplemented design option.
\end{abstract}
\clearpage
\tableofcontents
\clearpage

\section{Scope and design commitments}
\begin{enumerate}[leftmargin=*]
\item Preserve the corrected Seitz behavior, public identifiers, parameter
conventions, residuals, Jacobians and regression references. Permit only the
small internal hooks needed to reuse its assembly; verify unchanged Seitz results.
\item Add a distinct Mlika frictional formulation. Preserve the existing Mlika
frictionless energy formulation and its factory identifiers.
\item Use the existing two-dimensional line-facet ray mapping and the contact
time-interpolation configurations already supported by that mapping. The first
qualification uses the existing constant-in-time contact fixtures and Q1/Q2
spatial examples; higher temporal order and three-dimensional contact are not
asserted by this proposal. Retain the mapping component's current requirement
for compatible leader/follower basis counts; validate unsupported mixed-order
pairs explicitly. The total-Lagrangian reference geometry is fixed across the
trajectory, although a design perturbation changes it consistently at all steps.
\item Support $\theta\in\{-1,0,1\}$ with the same mathematical convention as the
thesis. MORIS labels $\theta=0$ ``neutral'' and $\theta=-1$ ``unsymmetric'';
the literature calls these non-symmetric and skew-symmetric, respectively.
Frictional tangents are generally non-symmetric, including for $\theta=1$.
\item Reconstruct previous configurations through previous-slab FI managers.
Do not use a traction-history map or add a friction ISV bank.
\item First target the existing elastic/hyperelastic constitutive models.
Constitutive internal variables for inelastic materials would require their own
accepted-state contract and are outside this proposal.
\end{enumerate}

\section{Literature target and distinction from the existing Mlika IWG}
The target is the large-deformation, ray-tracing Coulomb formulation of thesis
Sec.~4.2.1, Eq.~(4.25) \cite{mlika2018}, and paper Sec.~2.1, Eq.~(15)
\cite{mlika2017}. Its normal/friction operator is given by thesis
Eqs.~(4.17)--(4.21), and its generalized Newton tangent by Eq.~(4.32) and
Appendix~B. The earlier small-strain work \cite{chouly2018} treats a different
setting and does not replace this finite-deformation derivation.

The current MORIS Mlika normal residual instead follows the frictionless
potential-related family of thesis Eq.~(4.30), paper Eq.~(20). With $a_-=\min(a,0)$,
its directional contact residual is
\begin{equation}\label{eq:oldnormal}
 r_{\mathrm{energy}}(\delta\bm u)
 =\frac{1}{2\gamma}
 \left[-\theta t_n\,\delta t_n
       +a_-\bigl(\theta\delta t_n+\gamma\delta g\bigr)\right].
\end{equation}
The new implementation must not silently switch between this family at
$\mu=0$ and Eq.~(4.25) at $\mu>0$. It uses the vector formulation below for
all $\mu\geq0$, with its own frictionless limit and identifiers.

\section{Configuration, signs, frames, and discrete operators}
Consider one oriented contact quadrature point. Let $X$ be the fixed reference
leader point and $Y(\bm d)$ its current ray-mapped reference follower point.
The current and previous configurations are $\bm\varphi$ and
$\bm\varphi^-$, with $\bm\varphi(X)=X+\bm u(X)$. The ray equation is
\begin{equation}\label{eq:ray}
 \bm\varphi_\ell(X)-\bm\varphi_f(Y)+g\vn=\bm0.
\end{equation}
Here $g>0$ denotes separation, $\vn$ is the current outward leader unit normal,
and compression has $t_n\leq0$. The reference leader unit normal is $\vN$.
For the default PK1 constitutive mode,
\begin{equation}
 \vt=\bm\Pi\vN,\qquad t_n=\vt\cdot\vn,
 \qquad \Proj=\Id-\vn\vn\tr.
\end{equation}
$\bm\Pi$ denotes the first Piola--Kirchhoff stress. Both the normal and tangential
contact tractions are nominal forces per reference area, integrated using the
reference surface measure. A Cauchy traction must not be substituted without the
corresponding measure conversion.

Let $\bm d_\ell,\bm d_f$ denote current leader/follower coefficient vectors, and
$\Nm_\ell,\Nm_f(Y)$ the vector-valued FI shape matrices. Coefficients follow the
MORIS FI's field/component-major ordering, with temporal and spatial basis
ordering inherited from that FI. The contact geometry component converts its
internal node-major derivative ordering to this assembly ordering.

At fixed reference geometry and leader quadrature point, define
\begin{equation}\label{eq:T}
 \Tm_\ell=\frac{\partial\vt}{\partial\bm d_\ell},\qquad
 \frac{\partial\vt}{\partial\bm d_f}=\bm0.
\end{equation}
These are constitutive traction derivatives, not derivatives of the scalar
$t_n$; variation of the current normal is accounted for separately. The follower
shape matrix depends on current coefficients through $Y(\bm d)$.

\subsection{One oriented contribution and unbiased integration}
Every expression below is for one oriented quadrature contribution with factor
$1/2$. The unbiased formulation integrates over both potential contact surfaces.
The implementation audit must establish that MORIS's set/orientation construction
provides that coverage. Merely adding a factor $1/2$ to a single master/slave
pass is not an unbiased discretization. Quadrature, reference surface Jacobian,
selection, and thickness factors are applied exactly once.

\section{Previous-slab kinematics and the model time increment}
Thesis Eq.~(2.13) and the paper's objective backward-Euler sliding velocity
evaluate the previous configuration at the \emph{current} pairing $X,Y(\bm d)$:
\begin{equation}\label{eq:k}
 \bm k=\bm\varphi^-_\ell(X)-\bm\varphi^-_f(Y)+g\vn^-.
\end{equation}
With the actual positive model time/load increment $\Delta t$, define
\begin{equation}\label{eq:v}
 \bm v=-\frac{\bm k}{\Delta t},\qquad
 \vs=\Proj\bm k=-\Delta t\Proj\bm v.
\end{equation}
This includes previous normal rotation and follows the objective kinematics of
Refs.~\cite{mlika2018,mlika2017,poulios2015}. It is not a difference between
current and previous displacement values at independently remapped points.

The existing \code{GapData::mLeaderTangentVelocity} stores $\vs$, despite its
name: its construction contains no division by the time increment. Reuse that
increment and its derivatives, with
\begin{equation}\label{eq:beta}
 \beta_M=\beta=\frac{\gamma}{\Delta t},\qquad
 \vq_M=\Proj\vt+\beta\vs.
\end{equation}
The corrected Seitz IWG instead uses
$\vq_S=\Proj\vt+\gamma_{\tau,S}\vs$ \cite{seitz2019,seitz2019thesis}.
Seitz's paper explicitly distinguishes its slip increment from Mlika's velocity
immediately after Eq.~(34). No stress time-level difference remains.

\paragraph{Units and time convention.}
Use the time/load parameter actually supplied by the model. The earlier proposed
normalization $\tau=t/T_*$ and the \code{FrictionTimeScale} property are withdrawn:
there is no hidden or additional reference-time conversion. The initial strict
single-$\gamma$ qualification uses a consistently nondimensional problem with a
nondimensional model time/load parameter. This is an explicit implementation
scope, not a claim that the thesis prescribes a particular nondimensionalization.
It avoids treating a velocity and a length as dimensionally interchangeable.
Dimensional physical time would require an explicit problem nondimensionalization
or a separately derived normal/tangential parameter convention; the adapter must
not silently reinterpret physical seconds. Each benchmark must document its units,
model time grid and numerical $\gamma$. Changing the units or time parameterization
without consistently transforming the stabilization is not a meaningful comparison.

For example, a model increment $\Delta t=0.1$ gives $\beta=10\gamma$, not
$\beta=\gamma$. This is the velocity scaling in the selected discrete law; it
does not introduce a physical rate-dependent Coulomb friction model.

\subsection{Initialization and time-step selection}
The previous FI managers must represent the prescribed initial configuration at
the first increment and the accepted previous solution thereafter. The current
traction term below is present also at the first increment. Neither law has a
previous-traction initialization branch. Use actual current/previous evaluation
times or the time solver's step information to obtain $\Delta t$;
do not infer slab numbers by $\lfloor t/\Delta t\rfloor$. Variable step sizes
must be covered by tests. Require a positive $\Delta t$ and consistent previous
FI/geometry availability.

\section{Coulomb operator in a Seitz-like pointwise organization}
Define
\begin{align}
 a&=t_n+\gamma g,& m&=\min(a,0),& p&=-m\geq0,\\
 \rho&=\mu p=-\mu m,&
 \vq&=\Proj(\vt-\gamma\bm v)
      =\Proj\vt+\beta\vs.\label{eq:q}
\end{align}
The friction coefficient is $\mu\geq0$. Let
\begin{equation}\label{eq:projection}
 \bm z=\pi_\rho(\vq)=
 \begin{cases}
   \vq,&\|\vq\|\leq\rho,\\
   \displaystyle\rho\frac{\vq}{\|\vq\|},&\|\vq\|>\rho.
 \end{cases}
 \qquad
 \vC=m\vn+\bm z.
\end{equation}
For $\rho=0$, set $\bm z=\bm0$ without dividing by a norm. This handles
opening and frictionless contact. The normal compression and Coulomb radius
are based on the augmented argument $a$, not on $\gamma g$ alone.

This is the thesis operator $C_{\gamma,\mu}$, with the velocity convention in
Eq.~\eqref{eq:v}. Separate normal and tangential stabilization parameters would
be an additional formulation choice. The initial proposal retains the common
$\gamma$ of Eq.~(4.25).

\section{Weak form and residual blocks}
For test coefficients $\delta\bm d_\ell,\delta\bm d_f$, let
$\bm w_\ell=\Nm_\ell\delta\bm d_\ell$ and
$\bm w_f=\Nm_f(Y)\delta\bm d_f$. The oriented contact contribution in
thesis Eq.~(4.25) can be written
\begin{equation}\label{eq:weak}
 r_c(\bm w)=
 -\frac{\theta}{2\gamma}\vt\cdot\delta\vt
 +\frac{1}{2\gamma}\vC\cdot
      \left[\theta\delta\vt+\gamma(\bm w_f-\bm w_\ell)\right].
\end{equation}
Here $\delta\vt=\Tm_\ell\delta\bm d_\ell$ and the test jump means evaluation
of test functions at the current mapped point. It is not an instruction to add
an extra $\nabla\bm u_f\,\delta Y$ term to the residual test jump. Mapping
variation enters its \emph{linearization} through $\Nm_f(Y)$, as in the paper's
tangent system.

Let $w_q$ include the spatial quadrature weight, reference surface Jacobian,
and thickness, and put $\kappa=w_q/2$. The local residual increments are
\begin{align}
 \bm R_\ell&=-\kappa\Nm_\ell\tr\vC
 -\kappa\frac{\theta}{\gamma}\Tm_\ell\tr(\vt-\vC),\label{eq:Rl}\\
 \bm R_f&=\kappa\Nm_f\tr\vC.\label{eq:Rf}
\end{align}
These equations are the proposed implementation form. They follow Seitz's
leader/follower force and consistency-term organization. Both laws now use
current traction; only the slip-increment coefficient differs before parameter
matching.

For open contact, $\vC=\bm0$, but the leader stress-consistency term generally
remains. A test must distinguish zero physical contact force from a zero local
Nitsche residual at an arbitrary unconverged stress state.

\section{Semismooth derivatives of the contact operator}
For a current-state coefficient direction, assume initially that $\mu$,
$\Delta t$, $\gamma$, and $w_q$ are independent of displacement DOFs. Thickness
and selection data therefore have no displacement-DOF dependence in the initial
supported configuration. Geometry/design
dependence of $\gamma$ is addressed separately in Sec.~\ref{sec:design}.
Choose $\chi=1$ for $a<0$ and $\chi=0$ for $a>0$. Then
\begin{align}
 \delta a&=\vn\tr\delta\vt+\vt\tr\delta\vn+\gamma\delta g,\\
 \delta m&=\chi\delta a,& \delta\rho&=-\mu\delta m,\\
 \delta\Proj&=-\delta\vn\vn\tr-\vn\delta\vn\tr,\\
 \delta\vq&=\delta\Proj\,\vt+\Proj\delta\vt+\beta\delta\vs.
 \label{eq:dq}
\end{align}
For $r=\|\vq\|>0$ away from branch boundaries, define
\begin{equation}
 (\bm A,\bm b)=
 \begin{cases}
   (\Id,\bm0),&r<\rho,\\[3pt]
   \displaystyle\left(
     \frac{\rho}{r}\left(\Id-\frac{\vq\vq\tr}{r^2}\right),
     \frac{\vq}{r}\right),&r>\rho.
 \end{cases}
\end{equation}
The differential is
\begin{equation}\label{eq:dC}
 \delta\vC=\vn\delta m+m\delta\vn
            +\bm A\delta\vq+\bm b\delta\rho.
\end{equation}
The derivative of the tangential plane is already included in
Eq.~\eqref{eq:dq}; it must not be added a second time to the ball projection.

At $a=0$ or $r=\rho$, select and document a consistent one-sided generalized
Jacobian. An opening-side selection at $a=0$ and a specified stick/slip-side
selection at equality are suitable explicit conventions. Treat $(r,\rho)=(0,0)$
with a dedicated branch and no norm division. Branch locking may be used for
controlled finite-difference checks, not to freeze the physical branch across
Newton updates. Ordinary central FD agreement is required away from switching
surfaces; directional checks are needed on them.

\subsection{Geometric derivative content}
For current DOF variations, the previous coefficients remain fixed while the
current pairing changes. The raw velocity derivative is
\begin{equation}\label{eq:dv}
 \delta\bm v=\frac{1}{\Delta t}
 \left(\bm F^-_f(Y)\,\delta Y-\delta g\,\vn^-\right),
\end{equation}
matching thesis Eq.~(4.33) with the model time increment in Eq.~\eqref{eq:v}. The preferred MORIS
route is to use the existing derivatives of $\vs$ and the current normal in
\code{GapData}. These include the normal/projector and mapped follower
contributions. Both current and previous FI managers must be passed to
\code{Contact_Geometry}; the current Mlika frictionless calls pass null previous
managers and therefore cannot supply this friction kinematics.

\section{Current-state Jacobian blocks}
Let $b\in\{\ell,f\}$ and denote
$\vC_b=\partial\vC/\partial\bm d_b$ and
$\vt_b=\partial\vt/\partial\bm d_b$. Define the contracted derivatives
\begin{equation}
 \mathcal H_{\ell b}(\bm a)
 =\frac{\partial\Tm_\ell\tr}{\partial\bm d_b}\bm a,
 \qquad
 \mathcal M_{fb}(\bm a)
 =\frac{\partial\Nm_f\tr}{\partial\bm d_b}\bm a.
\end{equation}
The vector $\bm a$ is held fixed in these contractions; its derivative is
included in the separate product-rule term. For a standard leader
elastic/hyperelastic CM at fixed $X$, $\vt_f=0$ and
$\mathcal H_{\ell f}=0$. Differentiating Eqs.~\eqref{eq:Rl}--\eqref{eq:Rf} gives
\begin{align}
 \bm K_{\ell b}
 &=-\kappa\Nm_\ell\tr\vC_b
 -\kappa\frac{\theta}{\gamma}
 \left[\mathcal H_{\ell b}(\vt-\vC)
       +\Tm_\ell\tr(\vt_b-\vC_b)\right],\label{eq:Kl}\\
 \bm K_{fb}
 &=\kappa\left[\mathcal M_{fb}(\vC)+\Nm_f\tr\vC_b\right].\label{eq:Kf}
\end{align}
All four leader/follower blocks must be assembled. The follower test-function
mapping term $\mathcal M_{fb}$ is essential and is already represented by the
shared gap-data shape-derivative helpers.

The CM's traction derivative supplies $\Tm_\ell$; its contracted test-traction
derivative supplies $\mathcal H_{\ell\ell}$. For $\theta\ne0$, second derivatives
of constitutive traction are required. The old Mlika energy-form Hessian must
not be reused as if it were the Hessian of Eq.~\eqref{eq:weak}. Conversely, this
vector weak form does not require copying every second-gap/second-normal term
from that energy-form implementation.

\section{Previous-slab Jacobian}
Hold the current coefficients, current mapping, current gap, and current normal
fixed while perturbing previous coefficients. In the assumed elastic case,
$\vt$, $\Tm_\ell$, $m$, and $\rho$ then have no previous-state dependence.
Let
\begin{equation}
 \bm S^-_b=\frac{\partial\vs}{\partial\bm d^-_b}.
\end{equation}
In particular,
\begin{align}
 \bm S^-_\ell&=\Proj\left[\Nm^-_\ell+g
             \frac{\partial\vn^-}{\partial\bm d^-_\ell}\right],\\
 \bm S^-_f&=-\Proj\Nm^-_f(Y).
\end{align}
The previous follower field is evaluated at the current pair $Y$, not at a
separately recomputed previous contact pair. Thus
\begin{equation}
 \vq^-_b=\beta\bm S^-_b,\qquad
 \vC^-_b=\bm A\vq^-_b,
\end{equation}
and
\begin{align}
 \bm K^-_{\ell b}
 &=\kappa\left[-\Nm_\ell\tr
        +\frac{\theta}{\gamma}\Tm_\ell\tr\right]\vC^-_b,\label{eq:Kprevl}\\
 \bm K^-_{fb}&=\kappa\Nm_f\tr\vC^-_b.\label{eq:Kprevf}
\end{align}
Use \code{mLeaderdTangentVelocityduprev} and
\code{mLeaderdTangentVelocitydvprev}, with verified coefficient ordering.
This coupling needs a new \code{compute_jacobian_previous()} override.
It arises from previous displacement kinematics, not an independently stored
traction state. The corrected Seitz IWG has exactly this kinematic-only
previous-slab structure, with $\gamma_{\tau,S}$ replacing $\beta$. Its obsolete
previous-stress derivative has been removed. No constitutive-model switch to
previous FI managers belongs in the proposed Mlika implementation.

\section{Adjoints and design derivatives}\label{sec:design}
Let the assembled trajectory satisfy $\mathcal R_n=0$ and use the Lagrangian
\begin{equation}
 \mathcal L=J-\sum_n\bm\lambda_n\tr\mathcal R_n.
\end{equation}
The general discrete adjoint equations are
\begin{equation}\label{eq:adj}
 \sum_{m\geq n}\left(\frac{\partial\mathcal R_m}
                                {\partial\bm d_n}\right)\tr
          \bm\lambda_m
 =\left(\frac{\partial J}{\partial\bm d_n}\right)\tr.
\end{equation}
The new contact law contributes only the current Jacobian and the adjacent
previous-slab coupling, Eqs.~\eqref{eq:Kprevl}--\eqref{eq:Kprevf}. Existing DG
inertia/history terms can have longer coupling; their existing MORIS adjoint
handling remains in place. No contact ISV carry mechanism is added.

For a design variable $p$ and design-independent prescribed initial state,
\begin{equation}
 \frac{\dd J}{\dd p}=J_{,p}
 -\sum_n\bm\lambda_n\tr\mathcal R_{n,p}.
\end{equation}
If the initial state depends on $p$, its trajectory boundary terms must also be
included. State vectors are fixed while evaluating each partial
$\mathcal R_{n,p}$. Geometry sensitivities include reference normals, pairing,
shape functions, previous-field evaluations, integration measures, and
geometry-dependent stabilization. Property sensitivities include the CM and any explicitly design-dependent
input. In the initial input contract $\mu$ is a fixed scalar and the time grid
is design independent. If friction or time-grid quantities are later exposed as
ADVs, their explicit partial derivatives must be added to this same path. In
particular, $\delta\beta=\delta\gamma/\Delta t-
\gamma\,\delta(\Delta t)/(\Delta t)^2$ must then be included.

The first implementation uses the established MORIS geometry/material FD
sensitivity hooks for these partial design derivatives, with the analytic
current and previous Jacobians above. Every perturbation must invalidate the
contact geometry cache and reevaluate the previous fields at the perturbed
current pairing. No traction cache may survive a design perturbation.
Do not claim an analytic shape derivative merely because the state Jacobian is
analytic. Use the same end-to-end adjoint-versus-FD checks as for Seitz, including
multiple objective RHS columns. The accepted time grid is held fixed with
respect to the design variables in the initial qualification.

\section{Comparison with the corrected Seitz implementation}
The comparator is the corrected MORIS Seitz IWG \cite{moriscode}, checked against
Seitz's paper Eqs.~(33)--(35) and thesis Eqs.~(5.11)--(5.13)
\cite{seitz2019,seitz2019thesis}. Both the projection and the stress-consistency
term use \emph{current} constitutive traction. Previous configurations supply
slip kinematics only. The old comparison based on reconstructed previous
constitutive traction is superseded.

\begingroup\small
\setlength{\tabcolsep}{4pt}
\begin{longtable}{>{\raggedright\arraybackslash}p{0.20\textwidth}
>{\raggedright\arraybackslash}p{0.36\textwidth}
>{\raggedright\arraybackslash}p{0.36\textwidth}}
\toprule
Aspect & Corrected MORIS Seitz & Proposed Mlika Eq.~(4.25)\\
\midrule\endhead
Geometry & Shared current ray map and previous FI managers.
 & Same component and cache/invalidation contract.\\[3pt]
Tangential kinematics & Slip increment $\vs$ from shared gap data.
 & Same $\vs$, converted using $\beta=\gamma/\Delta t$.\\[3pt]
Trial vector & $\vq_S=\Proj\vt+\gamma_{\tau,S}\vs$.
 & $\vq_M=\Proj\vt+(\gamma/\Delta t)\vs$.\\[3pt]
Traction source & Current constitutive traction and its tangent.
 & Same; previous FIs supply only kinematics.\\[3pt]
Normal bound & $m=\min(t_n+\gamma g,0)$, $\rho=-\mu m$.
 & Same common-$\gamma$ normal argument and radius.\\[3pt]
Consistency terms & Current normal/tangential stress terms; the present MORIS code uses $\gamma$ in both denominators.
 & Same vector residual, Eqs.~\eqref{eq:Rl}--\eqref{eq:Rf}.\\[3pt]
Previous derivative & Slip kinematics only, scaled by $\gamma_{\tau,S}$.
 & Same kinematics, scaled by $\gamma/\Delta t$.\\[3pt]
Independent friction state & None; previous-stress helpers were removed.
 & None. No traction map or contact ISV update.\\[3pt]
Parameters & Existing four-scalar interface; explicit increment stiffness or its normal-SP fallback.
 & Separate two-scalar interface; velocity coefficient follows from $\gamma$ and $\Delta t$.\\
\bottomrule
\end{longtable}
\endgroup

\subsection{Matched-coefficient equivalence}
For the same smooth mapping branch, reference stress measure, $\theta$, $\mu$,
quadrature, selection/thickness factors and common normal/consistency parameter
$\gamma$, choose
\begin{equation}\label{eq:match}
 \gamma_{\tau,S}=\frac{\gamma}{\Delta t}.
\end{equation}
Then $\vq_S=\vq_M$, so both oriented residuals, all four current Jacobian
blocks and all four previous-slab blocks must agree. This is a stronger check
than the earlier proposal's frictionless comparison. At $\Delta t=1$, Seitz's
normal-SP fallback also gives this match. With varying $\Delta t$ or spatially
varying $\gamma$, matching must hold at each evaluation point; a single constant
Seitz parameter cannot generally provide that match throughout a trajectory.
Ordinary comparisons using the same numeric stiffness without this conversion
can differ at a finite mesh even though the continuum Coulomb law is the same.

In active contact, decomposing $\vt=t_n\vn+\Proj\vt$ gives
\begin{align}
 \bm R_\ell={}&-\kappa\Nm_\ell\tr(m\vn+\bm z)
 +\kappa\theta g\Tm_\ell\tr\vn
 -\kappa\frac{\theta}{\gamma}\Tm_\ell\tr(\Proj\vt-\bm z),\\
 \bm R_f={}&\kappa\Nm_f\tr(m\vn+\bm z).
\end{align}
These are the corrected Seitz normal/tangential blocks, with the trial coefficient
specified by the selected formulation. In particular, the shared assembly must
keep the stress-consistency denominator equal to $\gamma$; it must not substitute
$\beta$ there merely because $\beta$ multiplies the slip increment.

For $\mu=0$, both have $\bm z=0$. With matching full Nitsche settings, their
current residuals/tangents agree and their previous-contact Jacobians vanish.
This is the vector formulation's frictionless limit, not the older Mlika energy
form in Eq.~\eqref{eq:oldnormal}.

\subsection{What the corrected baseline establishes}
The independent two-block sticking patch changes bulk shear stress at exactly
zero interfacial slip. The corrected IWG transfers the analytic current traction
for different previous stresses, including reversal. Focused contact tests passed
610 assertions in 22 cases on each backend, including current and previous-slab
finite differences. Four Armadillo indenter forward/sensitivity cases passed with
unchanged references and tolerances. These checks support this corrected baseline;
they do not qualify a Mlika implementation, a full Eigen sensitivity workflow,
a three-dimensional mapping, or a complete global shear loading-cycle benchmark.

\section{Proposed MORIS structure and configuration}
\subsection{A small adapter over the corrected Seitz assembly}
Add \code{IWG_Isotropic_Struc_Nonlinear_Contact_Mlika_Frictional} with distinct
factory entries for all three $\theta$ variants. Preserve existing Mlika
energy-form identifiers and the corrected public Seitz behavior. Avoid a second
copy of the residual/Jacobian implementation.

The recommended arrangement is a thin derived adapter over the corrected Seitz
IWG, using a small protected virtual hook such as
\code{tangential_trial_coefficient(aGamma)}. The base implementation returns
Seitz's existing explicit increment stiffness or its SP fallback, preserving its
four-parameter interface. The Mlika override validates model time data and returns
$aGamma/\Delta t$. Replace the three local coefficient selections in residual,
current-Jacobian and previous-Jacobian evaluation by that hook. Do not repurpose
Seitz's second user parameter or mutate its parameter vector during evaluation.
The Mlika constructor assigns its own factory/type identity; it remains distinct
from the older Mlika energy formulation.

The shared assembly uses:
\begin{itemize}[leftmargin=*]
\item \code{Contact_Geometry} and \code{GapData} for current mapping, normals,
slip increments, and derivatives;
\item leader/follower current and previous FI managers, with time continuity enabled;
\item the existing \code{ElastLinIso} CM slot and the same selected stress measure;
\item \code{NitscheInterface} for the common $\gamma$ and existing thickness/select data;
\item Set residual/Jacobian assembly maps with unchanged field and side ordering.
\end{itemize}
All three methods must obtain the same coefficient from unchanged accepted
trajectory/time data. For fixed time data and DOF-independent $\gamma$, the hook
has no extra state derivative: both time-level Jacobians simply use its returned
coefficient. The $\theta/\gamma$ consistency prefactor is unchanged.

\subsection{One consistent projection evaluation}
The existing residual and tangent branches must be audited together before reuse.
In particular, zero Coulomb radius and zero trial norm require an explicit branch,
not division by zero; residual and both tangents must choose the same generalized
derivative at switching points. If needed, put that point-law evaluation in a
small protected helper shared by the assembly methods. Preserve gap-data branch
locking for controlled FD checks, not as a persistent friction state. These are
bounded internal changes with dedicated Seitz regressions, not a general rewrite
of the contact geometry or assembly.

Reuse the corrected analytic current and previous-slab Jacobians, retaining
constitutive test-traction and follower mapping derivatives. Combined assembly
needs an explicit override: the existing Seitz method
\code{compute_jacobian_and_residual()} currently raises an error. The Mlika
adapter must invoke the supported shared residual and Jacobian operations with
the same additive assembly and branch contract. Do not assume that inheriting this entry point supplies combined
assembly automatically. Reuse ordinary
MORIS partial-design FD dispatch deliberately; the shared direct
\code{compute_dRdp()} stub must not be advertised as an analytic or supported
standalone design operation. Audit the workflow dispatch and route a directly
exposed entry point to the supported FD path, or give a clear configuration error.

\subsection{Concrete initial input contract}
The new adapter's \code{set_parameters()} accepts exactly two scalar entries:
\begin{enumerate}[leftmargin=*]
\item friction coefficient $\mu\geq0$;
\item debug flag, zero for normal operation.
\end{enumerate}
It translates these once into the base class's full-Nitsche configuration, with
penalty-only disabled. Its coefficient hook overrides the base increment
stiffness choice. There is no \code{FrictionTimeScale} property, independent
friction history, separate tangential user stiffness, or penalty-only option in
the strict Eq.~(4.25) adapter. Seitz retains its existing options.

The SP supplies common positive $\gamma$. Validate finite inputs, $\mu\geq0$,
$\gamma>0$ and $\Delta t>0$, and consistent current/previous time metadata and FI
availability. The first qualification assumes nondimensional model time as
specified above, fixed accepted time grids, and DOF-independent $\mu$ and $\gamma$.
Geometry/material design dependence follows the partial-sensitivity path.
Unsupported state-dependent stabilization must fail explicitly until its full
product-rule derivatives are implemented. Higher temporal interpolation is not
implicitly qualified by this adapter.

\subsection{Evaluation sequence}
\begin{lstlisting}
configure and validate current/previous FI managers and accepted time data;
apply selection; obtain quadrature/thickness weight once;
map the leader point and evaluate GapData, including previous configurations;
follow the existing rejected-intersection policy if mapping fails;
evaluate current constitutive traction and requested derivatives;
obtain gamma; coefficient hook returns Seitz stiffness or Mlika gamma/DeltaT;
compute a, m, rho, q, z, C with one consistent branch description;
assemble requested residual, current Jacobian and/or previous Jacobian blocks;
leave forward fields and accepted solver trajectory unchanged.
\end{lstlisting}
Use local point results and allocate only requested derivative workspaces.
Invalidate geometry data on field/design changes. Combined assembly must use
consistent branch/state data and agree with separate calls. No evaluation writes
traction history, changes FI managers to reconstruct old stress, or copies the
removed Mlika penalty-friction prototype or its diagnostic output.

\section{Verification plan before qualification}
\begin{enumerate}[leftmargin=*]
\item \textbf{Point law.} Check opening, frictionless contact, stick, slip,
zero trial norm, zero radius, and branch-boundary directional derivatives.
Verify $\|\bm z\|\leq\rho$ and the converged physical friction law/dissipation.
\item \textbf{Element state derivatives.} Check every leader/follower current
Jacobian block against FD away from branch changes, for all $\theta$ variants.
Use nonmatching meshes and nontrivial normals and follower-coordinate changes.
Keep smooth FD comparisons away from contact-facet changes, and test facet
crossings separately as a robustness case.
\item \textbf{Previous-state derivatives.} Perturb previous leader and follower
coefficients independently, holding the current pairing fixed as required by the
partial derivative. Cover the first increment and unequal step sizes.
\item \textbf{Kinematics.} Check rigid translation, rigid rotation, sliding,
and reference/current traction-frame consistency. Verify the model time
convention and the explicit $1/\Delta t$ conversion. Include unequal increments.
\item \textbf{Independent sticking patch.} Reuse the analytic two-block patch:
zero slip, changing current bulk shear stress, and different previous stresses,
including reversal. Mlika must transfer the current analytic traction for several
positive $\Delta t$ values without acquiring an old-stress dependency.
\item \textbf{Matched Seitz comparison.} Set the Seitz increment coefficient
according to Eq.~\eqref{eq:match}. Compare residuals and every current/previous
Jacobian block for opening, stick and slip, all $\theta$ variants, and several
increments including $\Delta t\ne1$. Also check $\mu=0$ and zero-radius branches.
Use independent FD and patch references in addition to this shared-engine check.
\item \textbf{Assembly/lifecycle.} Check nonzero addends, repeated evaluation,
combined versus separate assembly, cache invalidation after geometry and field
changes, and absence of persistent friction-state writes.
\item \textbf{Forward and adjoint examples.} First reuse the parabolic indenter
geometry/loading in a separate Mlika driver configuration that leaves the
existing Seitz cases unchanged. Add the precompressed two-block shear loading
cycle for deliberate sticking, sliding, unloading and reversal. Add a multislab
adjoint check for a geometric ADV and a material/property ADV. Use independent
FD sweeps and the same physical configuration across solver/backend variants.
\item \textbf{Literature benchmark.} Reproduce the paper's elastic half-ring
case with and without friction, after the elementary tests pass. Document
material law, geometry, loading, surface integration, units and the model time grid.
\item \textbf{Backends and parallel runs.} Run the same tests in separate
Armadillo and Eigen builds and at appropriate MPI process counts. Use stored
physical references after qualification; treat timings and iteration counts as
diagnostics. Retain the corrected Seitz references/tolerances and add shared-engine
regressions that demonstrate its default coefficient policy is unchanged.
\end{enumerate}

\section{Review checkpoints}
The proposal remains tied to Eq.~(4.25) and now reuses the corrected Seitz
assembly. Before implementation, review the thin coefficient-policy adapter,
the nondimensional model-time scope, explicit $\gamma/\Delta t$ conversion,
matching conditions in Eq.~\eqref{eq:match}, unbiased integration coverage,
zero-radius branches, and supported initial spatial/time-interpolation scope.
The previous-traction baseline and the earlier additional time-scale property
are withdrawn. No Mlika code has been implemented by this document update.
Any internal Seitz hooks must preserve its corrected numerical behavior and
existing public interface; independent friction history is not proposed.

\begin{thebibliography}{9}
\bibitem{mlika2018}
R.~Mlika, \emph{Nitsche method for frictional contact and self-contact:
Mathematical and numerical study}, PhD thesis, Universit\'e de Lyon / INSA Lyon,
2018, NNT 2018LYSEI002. In particular Chapter~4, Eqs.~(4.17)--(4.25),
(4.30)--(4.33), and Appendix~B.
\url{https://theses.hal.science/tel-02067118/document}.

\bibitem{mlika2017}
R.~Mlika, Y.~Renard and F.~Chouly,
``An unbiased Nitsche's formulation of large deformation frictional contact and
self-contact,'' \emph{Computer Methods in Applied Mechanics and Engineering}
325 (2017), 265--288, DOI: \href{https://doi.org/10.1016/j.cma.2017.07.015}
{10.1016/j.cma.2017.07.015}.
Author manuscript: \url{https://hal.science/hal-01427872/document}.

\bibitem{seitz2019}
A.~Seitz, W.~A.~Wall and A.~Popp,
``Nitsche's method for finite deformation thermomechanical contact problems,''
\emph{Computational Mechanics} 63 (2019), 1091--1110,
DOI: \href{https://doi.org/10.1007/s00466-018-1638-x}{10.1007/s00466-018-1638-x}.
In particular Eqs.~(16), (33)--(35) and the paragraph after Eq.~(34).
The full paper supplied as \code{Seitz 2019a.pdf} was inspected.

\bibitem{seitz2019thesis}
A.~Seitz, \emph{Computational Methods for Thermo-Elasto-Plastic Contact},
PhD thesis, Technical University of Munich, accepted 2019.
In particular Sec.~5.1.1, Eqs.~(5.9)--(5.13), printed pp.~136--137.
\url{https://mediatum.ub.tum.de/doc/1454145/1454145.pdf}.

\bibitem{chouly2018}
F.~Chouly, R.~Mlika and Y.~Renard,
``An unbiased Nitsche's approximation of the frictional contact between two
elastic structures,'' \emph{Numerische Mathematik} 139 (2018), 593--631,
DOI: \href{https://doi.org/10.1007/s00211-018-0950-x}{10.1007/s00211-018-0950-x}.
\url{https://hal.science/hal-01240068/document}.

\bibitem{poulios2015}
K.~Poulios and Y.~Renard,
``An unconstrained integral approximation of large sliding frictional contact
between deformable solids,'' \emph{Computers \& Structures} 153 (2015), 75--90,
DOI: \href{https://doi.org/10.1016/j.compstruc.2015.02.027}
{10.1016/j.compstruc.2015.02.027}. Background on ray mapping and objective sliding.

\bibitem{moriscode}
{\raggedright
MORIS, branch \code{integration/contact-cmake-refactor}, working tree inspected
November 25, 2026. Relevant source files under \code{projects/FEM/INT/src/IWG/}:\par
\code{cl_FEM_IWG_Isotropic_Struc_Nonlinear_Contact_Seitz.cpp}\par
\code{cl_FEM_IWG_Isotropic_Struc_Nonlinear_Contact_Mlika.cpp}\par
\code{cl_FEM_Contact_Geometry.cpp}\par
\code{cl_FEM_Contact_Gap_Data.hpp}.\par}

The Seitz comparison uses the corrected current-traction implementation after
removal of the temporary alternative IWG. Spatial ray mapping, PK1 traction and
one-sided integration conventions remain MORIS adaptations. Literature equations
and the implementation comparator are identified separately throughout.
\end{thebibliography}
\end{document}
